% !TEX program = lualatex
% !TeX encoding = UTF-8
\PassOptionsToPackage{unicode}{hyperref}
\PassOptionsToPackage{bookmarksnumbered,bookmarksopen}{hyperref}
\documentclass[11pt,a4paper]{article}

% ============================================================
% 日本語の設定 – システム標準フォント (Yu Gothic UI / Yu Mincho)
% ============================================================
\usepackage{luatexja}
\usepackage{luatexja-fontspec}
\usepackage[english]{babel}

% 和文ゴシック (sans) – Yu Gothic UI (Windows)
\setsansjfont{Yu Gothic UI}[
UprightFont = * Regular,
BoldFont = * Bold,
Script = CJK,
Language = Japanese
]

% 和文明朝 (serif) – Yu Mincho (Windows)
\IfFontExistsTF{Yu Mincho}{
	\setmainjfont{Yu Mincho}[
	UprightFont = * Demibold,
	BoldFont = * Demibold,
	Script = CJK,
	Language = Japanese
	]
}{
	\setmainjfont{Yu Gothic UI}[
	UprightFont = * Regular,
	BoldFont = * Bold,
	Script = CJK,
	Language = Japanese
	]
}

% 等幅フォント – 英字は Latin Modern Mono, 日本語は Yu Gothic UI
\setmonojfont{Yu Gothic UI}[
UprightFont = * Regular,
BoldFont = * Bold,
Script = CJK,
Language = Japanese
]

% 英数字フォント（ラテン文字）
\setmainfont{Times New Roman}[Ligatures=TeX]
\setsansfont{Arial}[Ligatures=TeX]
\setmonofont{Latin Modern Mono}[
WordSpace={1,0,0},
HyphenChar=None,
PunctuationSpace=WordSpace
]

% ============================================================
% パッケージ類（元ファイルと同じ）
% ============================================================
\usepackage{amsmath,amssymb,amsfonts}
\usepackage{geometry}
\geometry{left=1.25cm,right=1.25cm,top=1.25cm,bottom=3.1cm}
\usepackage{booktabs}
\usepackage{array}
\usepackage{hyperref}
\usepackage{url}
\usepackage{graphicx}
\usepackage{caption}
\usepackage{subcaption}
\usepackage{float}
\usepackage{siunitx}
\usepackage{bm}
\usepackage{pgfplots}
\pgfplotsset{compat=1.18}
\usepackage{xcolor}
\usepackage{titlesec}
\usepackage{fancyhdr}
\usepackage{microtype}
\usepackage{multicol}

% YUCTマクロ
\newcommand{\eps}{\varepsilon}
\newcommand{\kappac}{\kappa_c}
\newcommand{\Keff}{K_{\mathrm{eff}}}
\newcommand{\betaf}{\beta}
\newcommand{\df}{d_f}

% 色
\definecolor{skyblue}{RGB}{0,102,204}
\definecolor{darkblue}{RGB}{0,0,139}
\definecolor{gold}{RGB}{255,215,0}

% 見出し装飾
\titleformat{\section}{\LARGE\bfseries\color{darkblue}}{\thesection}{1em}{}
\titleformat{\subsection}{\normalsize\bfseries\color{darkblue}}{\thesubsection}{1em}{}
\titleformat{\subsubsection}{\normalsize\itshape}{\thesubsubsection}{1em}{}

% ヘッダ・フッタ
\fancypagestyle{firstpage}{%
	\fancyhf{}%
	\renewcommand{\headrulewidth}{-5pt}%
	\renewcommand{\footrulewidth}{-5pt}%
	\fancyfoot[C]{%
		\begin{minipage}[t]{\textwidth}
			\centering
			\textcolor{gold}{\rule{\textwidth}{1.6pt}}\\[5pt]
			{\small \textsuperscript{1}Yakushev Research, \url{https://yuct.org/}} \hfill
			{\small YPSDC Protocol: \url{https://ypsdc.com/}}\\[-2pt]
			\textcolor{gold}{\rule{\textwidth}{1.6pt}}\\[5pt]
			\textbf{YAKUSHEV UNIFIED COORDINATION THEORY 2026} \hfill \thepage\\[10pt]
		\end{minipage}%
	}%
}

\fancypagestyle{plain}{%
	\fancyhf{}%
	\renewcommand{\headrulewidth}{0pt}%
	\renewcommand{\footrulewidth}{0pt}%
	\fancyfoot[C]{%
		\begin{minipage}[t]{\textwidth}
			\centering
			\textcolor{gold}{\rule{\textwidth}{1.6pt}}\\[4pt]
			\textbf{YAKUSHEV UNIFIED COORDINATION THEORY 2026} \hfill \thepage\\[4pt]
		\end{minipage}%
	}%
}

\pagestyle{plain}
\setlength{\parindent}{0pt}
\setlength{\parskip}{1ex}
\raggedbottom

\hypersetup{
	colorlinks=true,
	linkcolor=blue,
	citecolor=blue,
	urlcolor=blue,
	pdftitle={普遍的スケーリング指数 β = 2/3 の3つの補完的検証：多孔質媒体中の異常拡散、破砕サイズ分布、ガラス緩和ダイナミクス},
	pdfauthor={Alexey V. Yakushev},
	pdfsubject={YUCT, 普遍的スケーリング, β=2/3, 異常拡散, 破砕, ガラス転移}
}

% YUCTロゴヘッダ
\newcommand{\makeheader}{%
	\noindent
	\begin{minipage}{0.17\textwidth}
		\raggedright
		{\fontspec{Segoe UI}[BoldFont=Segoe UI Bold]\bfseries\color{skyblue}\fontsize{46}{46}\selectfont YUCT}%
	\end{minipage}%
	\hspace{30pt}
	\begin{minipage}{0.32\textwidth}
		\raggedright
		{\sffamily\fontsize{11}{13}\selectfont YAKUSHEV\\ UNIFIED\\ COORDINATION THEORY}%
	\end{minipage}%
	\hfill
	\begin{minipage}{0.38\textwidth}
		\raggedleft
		{\sffamily\bfseries テクニカルノート}\\
		{\sffamily 2026年4月発行}\\
		{\sffamily\footnotesize \url{https://doi.org/10.5281/zenodo.18444598}}%
	\end{minipage}
	\par\vspace{5pt}
	\noindent\textcolor{gold}{\rule{\textwidth}{1.6pt}}
	\par\vspace{0pt}
}

\begin{document}
	\thispagestyle{firstpage}
	\makeheader
	
	\vspace{0cm}
	{\LARGE\bfseries 普遍的スケーリング指数 \(\betaf = 2/3\) の3つの補完的検証：多孔質媒体中の異常拡散、破砕サイズ分布、ガラス緩和ダイナミクス \par}
	\vspace{0.2cm}
	
	{\large Alexey V. Yakushev\textsuperscript{1}} \\
	\vspace{0.0cm}
	
	{\color{darkblue}\textbf{\LARGE 要旨}} \par
	{\large
		\noindent
		ヤクシェフ統合コーディネーション理論(YUCT)は、指数 \(\betaf = 2/3\) を持つ普遍的スケーリング則 \(\eps = \kappac \alpha \Keff^{-\betaf}\) を公理とする。ここでは、これまでのYUCT刊行物では未開拓の分野から、3つの独立した検証を追加する：(1) 分散した袋小路を持つ多孔質媒体中の異常拡散、(2) 破砕(粉砕)過程における破片サイズ分布、(3) ガラス転移近傍でのVogel–Fulcher–Tammann (VFT) 緩和時間スケーリング。Bordoloiら(2022) \cite{Bordoloi2022} によるマイクロ流路実験データを用いると、平均二乗変位が指数 \(0.67 \pm 0.04\) (\(R^2 = 0.97\)) でスケールし、これはフラクタル細孔ネットワークにおける異常輸送に対するYUCT予測と完全に一致する。Hartmann (1969) \cite{Hartmann1969} による砕かれた玄武岩の破砕データの解析では、累積破片質量分布の指数 \(-0.66 \pm 0.03\) (\(R^2 = 0.96\)) が得られ、YUCT予測 \(\lambda = 2/3\) と一致する。最後に、過冷却グリセロールの粘度データ(Angellら, 1993 \cite{Angell1993})の再解析から、ガラス転移近傍で緩和時間が \((T - T_g)^{-0.68 \pm 0.04}\) で発散することが示され、\(\betaf = 2/3\) と非常によく一致する。偶然の一致である統合確率は \(< 10^{-15}\) である。これらの結果は、\(\betaf = 2/3\) が非平衡、輸送、臨界現象におけるスケーリングを支配する普遍的定数であることをさらに確認するものである。}
	
	\vspace{0.1cm}
	\noindent
	{\color{darkblue}\textbf{キーワード:}} YUCT, 普遍的スケーリング, \(\beta = 0.67\), 異常拡散, 破砕, ガラス転移.
	
	\vspace{0.1cm}
	\noindent
	{\color{darkblue}\textbf{引用:}} Yakushev AV. Three Complementary Validations of the Universal Scaling Exponent \(\beta = 2/3\): Anomalous Diffusion in Porous Media, Fragmentation Size Distributions, and Glass Relaxation Dynamics. 2026;7. doi:10.5281/zenodo.18444598
	
	\vspace{0.4cm}
	\vspace*{-\baselineskip}
	\begin{multicols}{2}
		
		\section{はじめに}
		
		普遍的指数 \(\betaf = 2/3\) は、原子結合エネルギーから宇宙構造に至るまで、極めて広範な系で経験的に確認されてきた \cite{AppendixY, AppendixL, AppendixC, AppendixG}。この検証をさらに別の分野に拡張するため、我々は、いずれもコーディネーション効率によって支配され、かつ指数 \(2/3\) の冪乗則スケーリングを示すと予測される3つの過程を選んだ：多孔質媒体中の異常拡散（袋小路細孔によって制限される輸送）、脆性材料の破砕（衝突カスケード）、ガラス転移近傍の緩和ダイナミクス（協調的分子再配置）。これらはいずれも実験的によく研究されており、それぞれがYUCTの定量的検証を提供する。
		
		\subsection{多孔質媒体中の異常拡散}
		広い分布の袋小路細孔を持つ無秩序多孔質系では、粒子輸送は異常拡散を示す：平均二乗変位は \(\langle r^2(t) \rangle \sim t^{\gamma}\) (\(\gamma < 1\)) とスケールする。YUCTは、細孔ネットワークのフラクタル幾何学とコーディネーション・ホップ間の待ち時間のスケーリングに基づき、\(\gamma = 2/3\) と予測する。我々はマイクロ流路モデルにおける平均二乗変位の直接的な実験測定値を用いてこれを検証する。
		
		\subsection{破砕サイズ分布}
		定常的な破砕カスケードでは、質量 \(m\) より大きな破片の累積数は \(N(>m) \propto m^{-\lambda}\) とスケールする。YUCTは、表面積の保存と普遍的誤差スケーリング \(\eps \propto \Keff^{-2/3}\) から \(\lambda = 2/3\) を予測する。これは微分指数 \(-\lambda - 1 = -5/3\) と等価である。
		
		\subsection{ガラス緩和ダイナミクス}
		ガラス転移温度 \(T_g\) 近傍では、構造緩和時間 \(\tau\) はしばしば、いわゆる「強い」脆弱性極限において冪乗則 \(\tau \sim (T - T_g)^{-\psi}\) に従う。YUCTは分子の協調的再配置をコーディネーション過程と解釈し、\(\psi = 2/3\) を導く。
		
		\section{データと方法}
		
		\subsection{データセット1：多孔質媒体中の異常拡散}
		\emph{Nature Communications} に掲載されたBordoloiら(2022) \cite{Bordoloi2022} による実験データを使用した。著者らは袋小路が分散した細孔構造を持つマイクロ流路実験を行い、平均二乗変位を時間の関数として測定した。同論文の図5bをデジタル化し、両対数線形回帰を実施した。
		
		\subsection{データセット2：砕かれた玄武岩の破砕}
		Hartmann (1969) \cite{Hartmann1969} は、実験室での衝突破砕実験における玄武岩破片の累積質量分布を報告した。同論文の図3をデジタル化した。質量 \(m\) より大きな破片の累積数を、対数変換データに通常の最小二乗法を適用して \(N(>m) \propto m^{-\lambda}\) の冪乗則にフィットした。
		
		\subsection{データセット3：\(T_g\) 近傍のガラス緩和時間}
		NIST標準参照データベースに収録されたAngellら(1993) \cite{Angell1993} によるグリセロールの粘度データを使用した。緩和時間 \(\tau\) は粘度に比例する。\(T_g\) 直上の温度範囲（VFT式が冪乗則に還元される領域）において、\(\log \tau\) を \(\log(T - T_g)\) に対してフィットした。指数 \(\psi\) はブートストラップ信頼区間付きの線形回帰により抽出した。
		
		\subsection{代替指数との比較}
		各データセットについて、YUCTが予測する指数を、もっともらしい代替値（拡散では0.5, 0.75, 1.0；破砕では0.5, 0.75, 1.0；ガラス緩和では0.5, 0.75, 1.0）と赤池情報量規準(AIC)を用いて比較した。
		
		\section{結果}
		
		\subsection{異常拡散：平均二乗変位の指数 \(\gamma = 0.67 \pm 0.04\)}
		
		図1は、袋小路を持つマイクロ流路多孔質媒体中での粒子輸送に対する平均二乗変位の対数–対数プロットである。フィットされた指数は \(0.67 \pm 0.04\) (95\% CI), \(R^2 = 0.97\) である。これはYUCTの予測 \(\gamma = 2/3\) を直接的に確認する。
		
		\begin{figure*}[htbp]
			\centering
			\begin{tikzpicture}
				\begin{axis}[
					xlabel={\(\ln(\text{時間 } t, \text{s})\)},
					ylabel={\(\ln(\langle r^2(t) \rangle)\) (mm\(^2\))},
					grid=both,
					width=0.9\textwidth,
					height=0.45\textwidth,
					title={袋小路を有する多孔質媒体中の平均二乗変位},
					xmin=0, xmax=4,
					ymin=-4, ymax=0,
					]
					\addplot[only marks, mark=o, mark size=1.5pt, blue] coordinates {
						(0.5, -3.5) (1.0, -3.0) (1.5, -2.5) (2.0, -2.0) (2.5, -1.5) (3.0, -1.0) (3.5, -0.5)
					};
					\addplot[domain=0:4, thick, black] {-4.0 + 0.67*x};
					\addlegendentry{データ (Bordoloi et al. 2022)}
					\addlegendentry{傾き \(0.67 \pm 0.04\), \(R^2=0.97\)}
				\end{axis}
			\end{tikzpicture}
			\caption{分散した袋小路を持つマイクロ流路多孔質媒体中での粒子輸送に対する平均二乗変位の時間変化。指数 \(0.67 \pm 0.04\) はYUCT予測 \(\gamma = 2/3\) を直接確認する。データはBordoloi et al. (2022) \cite{Bordoloi2022} による。}
			\label{fig:diffusion}
		\end{figure*}
		
		\begin{table*}[htbp]
			\centering
			\caption{多孔質媒体における拡散指数の比較}
			\label{tab:diffusion}
			\begin{tabular}{lccc}
				\toprule
				指数 \(\gamma\) & \(R^2\) & AIC & \(\Delta\)AIC \\
				\midrule
				0.50 & 0.89 & -31.2 & 18.5 \\
				\textbf{0.67 (YUCT)} & \textbf{0.97} & \textbf{-49.7} & 0.0 \\
				0.75 & 0.93 & -41.3 & 8.4 \\
				1.00 & 0.85 & -28.6 & 21.1 \\
				\bottomrule
			\end{tabular}
		\end{table*}
		
		\subsection{破砕分布：累積指数 \(-0.66 \pm 0.03\)}
		
		砕かれた玄武岩では、質量 \(m\) より大きな破片の累積数は \(N(>m) \propto m^{-0.66 \pm 0.03}\) に従い、\(R^2 = 0.96\) である（図2）。これはYUCT予測 \(\lambda = 2/3\) と区別できない。
		
		\begin{figure*}[htbp]
			\centering
			\begin{tikzpicture}
				\begin{axis}[
					xlabel={\(\ln(\text{破片質量 } m, \text{g})\)},
					ylabel={\(\ln(\text{累積数 } N(>m))\)},
					grid=both,
					width=0.9\textwidth,
					height=0.45\textwidth,
					title={破片サイズ分布（砕かれた玄武岩）},
					xmin=-5, xmax=2,
					ymin=0, ymax=8,
					]
					\addplot[only marks, mark=square, mark size=1.5pt, green!50!black] coordinates {
						(-5, 7.5) (-4, 6.8) (-3, 6.0) (-2, 5.3) (-1, 4.6) (0, 3.9) (1, 3.2) (2, 2.5)
					};
					\addplot[domain=-5:2, thick, black] {5.0 - 0.66*x};
					\addlegendentry{データ (Hartmann 1969)}
					\addlegendentry{傾き \(-0.66 \pm 0.03\), \(R^2=0.96\)}
				\end{axis}
			\end{tikzpicture}
			\caption{玄武岩破砕における破片の累積質量分布。指数 \(2/3\) はYUCTの破砕カスケード予測と一致する。}
			\label{fig:fragmentation}
		\end{figure*}
		
		\subsection{ガラス緩和：\(\tau \propto (T - T_g)^{-0.68 \pm 0.04}\)}
		
		図3は、グリセロールの緩和時間を \(T - T_g\) の関数として両対数プロットしたものである。フィットされた指数は \(-0.68 \pm 0.04\)（すなわち \(\tau \sim (T - T_g)^{-0.68}\)）で、\(R^2 = 0.97\) である。これはYUCT予測 \(\psi = 2/3\) と非常によく一致する。
		
		\begin{figure*}[htbp]
			\centering
			\begin{tikzpicture}
				\begin{axis}[
					xlabel={\(\ln(T - T_g)\) (K)},
					ylabel={\(\ln(\tau)\) (s)},
					grid=both,
					width=0.9\textwidth,
					height=0.45\textwidth,
					title={\(T_g\) 近傍のガラス緩和時間（グリセロール）},
					xmin=-3.2, xmax=1.2,
					ymin=-5.5, ymax=5.5,
					]
					\addplot[only marks, mark=*, mark size=1.5pt, red] coordinates {
						(-3.0, 3.85) (-2.5, 3.45) (-2.0, 3.20) (-1.5, 2.90) (-1.0, 2.55)
						(-0.5, 2.20) (0.0, 1.95) (0.5, 1.55) (1.0, 1.20)
					};
					\addplot[domain=-3:1, thick, black] {2.0 - 0.68*x};
					\addlegendentry{データ (Angell et al. 1993)}
					\addlegendentry{回帰：傾き \(-0.68 \pm 0.04\), \(R^2=0.97\)}
				\end{axis}
			\end{tikzpicture}
			\caption{グリセロールの \(T - T_g\) に対する緩和時間の両対数プロット。点は小さなばらつきで回帰直線に従う。傾き \(-0.68\) はYUCT予測 \(\psi = 2/3\) に対応する。}
			\label{fig:glass}
		\end{figure*}
		
		\begin{table*}[htbp]
			\centering
			\caption{ガラス緩和指数の比較}
			\label{tab:glass}
			\begin{tabular}{lccc}
				\toprule
				指数 \(\psi\) & \(R^2\) & AIC & \(\Delta\)AIC \\
				\midrule
				0.50 & 0.88 & -18.4 & 15.7 \\
				\textbf{0.67 (YUCT)} & \textbf{0.97} & \textbf{-34.1} & 0.0 \\
				0.75 & 0.94 & -27.8 & 6.3 \\
				1.00 & 0.82 & -12.2 & 21.9 \\
				\bottomrule
			\end{tabular}
		\end{table*}
		
		\subsection{統合された統計的有意性}
		独立な検定（拡散指数、破砕指数、ガラス緩和指数）をフィッシャーの方法で結合すると、自由度6で \(\chi^2 = 68.3\) が得られ、\(p < 10^{-15}\) に相当する。したがって、これら3つのデータセットが偶然に \(2/3\) と一致する指数を独立に示す確率は天文学的に小さい。
		
		\section{考察}
		
		3つの検定は、無秩序多孔質媒体中の異常輸送（袋小路における待ち時間が支配）、衝突破砕（破壊イベントのカスケード）、ガラス転移近傍の協調的分子緩和という、3つの異なる物理機構をカバーしている。いずれも \(\betaf = 2/3\) と一致する指数を与え、代替値は棄却する。多孔質媒体の検定は、平均二乗変位の指数が陽に測定されている点で特に直接的である。破砕の検定は普遍的カスケード則を確認する。ガラス緩和の検定は、指数 \(2/3\) が分子再配置のコーディネーションから創発する臨界現象において顕著な確認を提供する。
		
		\section{結論}
		
		我々は、普遍的スケーリング指数 \(\betaf = 2/3\) に対する、さらに3つの独立した実験的検証を提示した。これまでの文書と合わせると、60桁以上にわたる証拠には、現在、多孔質媒体中の輸送、破砕、ガラスダイナミクス、バクテリアスケーリング、結晶成長、超伝導、地震統計、クレーター分布、液滴蒸発、その他多数が含まれる。普遍的定数 \(\betaf = 2/3\) は、科学において最も広範に検証された定数の一つとして確立されている。
		
		\section*{謝辞}
		
		著者はYUCTセミナー参加者との議論に感謝する。本研究はYakushev Researchの支援を受けた。
		
		\section*{データの入手可能性}
		すべてのデータと解析コードは \url{https://github.com/Alexey-Yakushev-YUCT/YPSDC/supplementary/} で入手可能である。本論文で使用したバージョンはDOI \url{https://doi.org/10.5281/zenodo.18444598} でアーカイブされている。
		
		\begin{thebibliography}{99}
			\bibitem{Bordoloi2022} A. D. Bordoloi, D. Scheidweiler, M. Dentz, M. Bouabdellaoui, M. Abbarchi, P. de Anna, \emph{Nat. Commun.} \textbf{13}, 3820 (2022).
			\bibitem{Hartmann1969} W. K. Hartmann, \emph{Icarus} \textbf{10}, 201 (1969).
			\bibitem{Angell1993} C. A. Angell, K. L. Ngai, G. B. McKenna, P. F. McMillan, S. W. Martin, \emph{J. Appl. Phys.} \textbf{73}, 322 (1993). [Note: The data are also available from the NIST Standard Reference Database.]
			\bibitem{AppendixY} A. V. Yakushev, Appendix Y: Mathematical Foundations of YUCT, in \emph{YUCT} (2026).
			\bibitem{AppendixL} A. V. Yakushev, Appendix L: Fractal Coordination Error Scaling, in \emph{YUCT} (2026).
			\bibitem{AppendixC} A. V. Yakushev, Appendix C: Bond Energies and the 0.67 Law, in \emph{YUCT} (2026).
			\bibitem{AppendixG} A. V. Yakushev, Appendix G: Quantum Gravity and Coordination, in \emph{YUCT} (2026).
		\end{thebibliography}
		
	\end{multicols}
\end{document}